\chapter{Mathematical Derivations under Normal Priors}
\label{sec:appendix_normal_priors}

\section{Derivation of the Posterior Decision Rule}

\paragraph{Setup}
Let $\bar{y}_C$ and $\bar{y}_T$ denote the observed sample means for the control and treatment arms, respectively. Recall that we assume that lower values of $y$ are clinically beneficial.  

Conditional on the true underlying parameters, the sampling distributions of the arm means are given by:
\begin{align*}
    \bar{y}_C \mid \theta_C &\sim \mathcal{N}(\theta_C, v_C^2) \quad \text{where} \quad v_C^2 = \frac{\sigma_C^2}{n_C}  \\
    \bar{y}_T \mid \theta_T &\sim \mathcal{N}(\theta_T, v_T^2) \quad \text{where} \quad v_T^2 = \frac{\sigma_T^2}{n_T}, 
\end{align*}
where $\sigma_C^2$ and $\sigma_T^2$ are the known variances, and $n_C, n_T$ are the sample sizes. The true treatment effect is defined as $\delta = \theta_C - \theta_T$. 

During the design evaluation, the true parameter $\theta_C$ is either fixed at a specific scalar value $\theta_C^*$ (when computing conditional T1E rate or power) or distributed following the design prior $\pi_D(\theta_C)$ (when calculating $\avgTIE$ or $\avgPow$).


\paragraph{Analysis Prior and Posterior}
We incorporate an informative historical analysis prior for the control arm, $\pi_A(\theta_C) \sim \mathcal{N}(\mu_{A,C}, \sigma_{A,C}^2)$, where $\mu_{A,C}$ represents the historical mean and $\sigma_{A,C}^2$ quantifies the prior uncertainty. We assume a non-informative analysis prior for the treatment arm.
Due to Normal-Normal conjugacy, the posterior distribution for the control parameter remains Normally distributed, $\theta_C \mid \bar{y}_C \sim \mathcal{N}(\mu_{\text{post},C}, \sigma_{\text{post},C}^2)$, where:
$$
    \sigma_{\text{post},C}^2 = \left( \frac{1}{\sigma_{A,C}^2} + \frac{1}{v_C^2} \right)^{-1} = \frac{\sigma_{A,C}^2 v_C^2}{\sigma_{A,C}^2 + v_C^2},
$$
and: 
$$
    \mu_{\text{post},C} = (1 - W)\mu_{A,C} + W\bar{y}_C,
$$
given the data weight: 
$$
    W = \frac{1/v_C^2}{1/\sigma_{A,C}^2 + 1/v_C^2} = \frac{\sigma_{A,C}^2}{\sigma_{A,C}^2 + v_C^2},
$$
where $(1 - W)$ quantifies the weight of the historical prior on the final control estimate. 

Conversely, we assume a non-informative (vague) analysis prior for $\theta_T$. Thus, the posterior for the treatment arm is driven entirely by the likelihood:
$$
    \theta_T \mid \bar{y}_T \sim \mathcal{N}(\bar{y}_T, v_T^2).
$$

\paragraph{Decision Rule}
Assuming independence between the two trial arms, the posterior distribution of the treatment contrast $\delta = \theta_C - \theta_T$ is also Normal:
$$
    \delta \mid \bar{y}_C, \bar{y}_T \sim \mathcal{N}(\hat{\delta}, V_{\delta}),
$$
where the posterior mean and variance are exactly:
\begin{align*}
    \hat{\delta} &= \mu_{\text{post},C} - \bar{y}_T \\
    V_{\delta} &= \sigma_{\text{post},C}^2 + v_T^2.
\end{align*}

A trial is declared successful if the posterior probability that the treatment effect exceeds a critical threshold $q$ (e.g., $q=0$ for any benefit, or $q=\delta_0$ for clinical relevance) is greater than a pre-specified confidence level $p$:
\begin{align}
    \Pr(\delta > q \mid \bar{y}_C, \bar{y}_T) &> p \nonumber \\
    1 - \Phi\left( \frac{q - \hat{\delta}}{\sqrt{V_{\delta}}} \right) &> p. \label{eq:app_decision_rule}
\end{align}
If Inequality~\eqref{eq:app_decision_rule} holds, the null hypothesis is rejected.

We can reformulate this inequality directly in terms of $\hat{\delta}$:
\begin{proof}
Starting from the posterior decision rule established in Equation~\eqref{eq:app_decision_rule}:
$$
    1 - \Phi\left( \frac{q - \hat{\delta}}{\sqrt{V_{\delta}}} \right) > p \implies \Phi\left( \frac{q - \hat{\delta}}{\sqrt{V_{\delta}}} \right) < 1 - p \implies \frac{q - \hat{\delta}}{\sqrt{V_{\delta}}} < \Phi^{-1}(1 - p).
$$
Using the symmetry of the standard Normal distribution, $\Phi^{-1}(1 - p) = - \Phi^{-1}(p)$. Defining $z_p = \Phi^{-1}(p)$ as the critical $z$-value corresponding to the required posterior probability threshold $p$ (e.g., if $p = 0.975$, $z_p \approx 1.96$), we obtain:
$$
    \frac{q - \hat{\delta}}{\sqrt{V_{\delta}}} < -z_p.
$$
Multiplying both sides by the strictly positive posterior standard deviation $\sqrt{V_{\delta}}$ and rearranging to isolate $\hat{\delta}$, we arrive at the final condition for study success:
\begin{equation} \label{eq:app_success_condition_delta}
    \hat{\delta} > q + z_p \sqrt{V_{\delta}}.
\end{equation}
\end{proof}

We denote the deterministic right hand side of Equation~\eqref{eq:app_success_condition_delta} as $k$:
$$
    k = q + z_p \sqrt{V_{\delta}}.
$$
Because the posterior variance $V_{\delta}$ depends exclusively on the fixed prior variance $\sigma_{A,C}^2$ and the known sampling variances ($v_C^2, v_T^2$), the threshold $k$ is a constant determined entirely at the design stage.


\section{Conditional power and T1E rate}

To evaluate the operating characteristics under theoretical repetitions of the trial, we derive the frequentist sampling distribution of $\hat{\delta}$.

\paragraph{Expected Value of the Estimator}
By linearity of expectation, and substituting $\Erw[\bar{y}_C \mid \theta_C] = \theta_C$ and $\Erw[\bar{y}_T \mid \theta_C, \delta] = \theta_C - \delta$, we obtain:
\begin{align}
    \Erw[\hat{\delta} \mid \theta_C, \delta] &= \Erw[\mu_{\text{post},C} \mid \theta_C] - \Erw[\bar{y}_T \mid \theta_C, \delta] \nonumber \\
    &= (1 - W)\mu_{A,C} + W \Erw[\bar{y}_C \mid \theta_C] - (\theta_C - \delta) \nonumber \\
    &= (1 - W)\mu_{A,C} + W\theta_C - \theta_C + \delta \nonumber \\
    &= (1 - W)(\mu_{A,C} - \theta_C) + \delta. \label{eq:app_expected_delta}
\end{align}

\paragraph{Sampling Variance of the Estimator}
Assuming the control and treatment arms are enrolled independently:
\begin{align}
    \Var(\hat{\delta} \mid \theta_C, \delta) &= \Var(\mu_{\text{post},C} \mid \theta_C) + \Var(\bar{y}_T \mid \theta_C, \delta) \nonumber \\
    &= Var\Big((1 - W)\mu_{A,C} + W\bar{y}_C\Big) + \Var(\bar{y}_T) \nonumber \\
    &= W^2 \Var(\bar{y}_C) + v_T^2 \nonumber \\
    &= W^2 v_C^2 + v_T^2. \label{eq:app_variance_delta}
\end{align}
We denote this sampling variance as $V_{\text{samp}} = W^2 v_C^2 + v_T^2$. The sampling distribution is thus:
$$
    \hat{\delta} \mid \theta_C, \delta \sim \mathcal{N}\Big( (1 - W)(\mu_{A,C} - \theta_C) + \delta, \; V_{\text{samp}} \Big).
$$

\paragraph{Closed-Form Solutions}
To compute the conditional power, we replace $\theta_C$ with a true control rate $\theta_C^*$ and integrate the sampling distribution over the success region defined by $k$:
\begin{align*}
    \Pow(\theta_C^*, \delta) &= \Pr(\hat{\delta} > k \mid \theta_C^*, \delta) \\
    &= 1 - \Phi\left( \frac{k - \Erw[\hat{\delta} \mid \theta_C^*, \delta]}{\sqrt{V_{\text{samp}}}} \right). 
\end{align*}
Substituting the threshold $k$, we obtain the conditional power:
\begin{equation} \label{eq:app_power_closed_form}
    \Pow(\theta_C^*, \delta) = 1 - \Phi\left( \frac{q + z_p \sqrt{V_{\delta}} - \big[(1 - W)(\mu_{A,C} - \theta_C^*) + \delta\big]}{\sqrt{V_{\text{samp}}}} \right),
\end{equation}
which reduces to the conditional T1E rate by setting $\delta = 0$: 
\begin{equation} \label{eq:app_t1e_closed_form}
    \TIE(\theta_C^*) = 1 - \Phi\left( \frac{q + z_p \sqrt{V_{\delta}} - (1 - W)(\mu_{A,C} - \theta_C^*)}{\sqrt{V_{\text{samp}}}} \right). 
\end{equation}

with $V_{\delta} = \frac{\sigma_{A,C}^2 v_C^2}{\sigma_{A,C}^2 + v_C^2} + v_T^2$.


\section{Average power and T1E rate}
\label{sec:avgmetrics}
When the true control parameter is distributed according to a Normal design prior, $\theta_C \sim \pi_D(\theta_C) = \mathcal{N}(\mu_{D,C}, \sigma_{D,C}^2)$, we have to obtain the marginal sampling distribution of $\hat{\delta}$.

\paragraph{Marginal Expectation and Variance}
Using the Law of Total Expectation over $\pi_D(\theta_C)$:
    $$\Erw[\hat{\delta} \mid \delta] = E_{\theta_C} \Big[ \Erw[\hat{\delta} \mid \theta_C, \delta] \Big] $$

Substituting Equation~\eqref{eq:app_expected_delta}, we obtain:
\begin{equation} \label{eq:app_marginal_expectation}
    E_{\theta_C} \Big[ (1 - W)(\mu_{A,C} - \theta_C) + \delta \Big] = (1 - W)(\mu_{A,C} - \mu_{D,C}) + \delta. 
\end{equation}
Applying the Law of Total Variance:
\begin{align}
\Var(\hat{\delta} \mid \delta) &= E_{\theta_C} \Big[ \Var(\hat{\delta} \mid \theta_C, \delta) \Big] + \Var_{\theta_C} \Big( \Erw[\hat{\delta} \mid \theta_C, \delta] \Big) \nonumber \\
&= E_{\theta_C} \Big[ V_{\text{samp}} \Big] + \Var_{\theta_C} \Big( (1 - W)(\mu_{A,C} - \theta_C) + \delta \Big) \nonumber \\
&= V_{\text{samp}} + (1 - W)^2 \, \Var(\theta_C) \nonumber \\
&= W^2 v_C^2 + v_T^2 + (1 - W)^2 \sigma_{D,C}^2. \label{eq:app_marginal_variance}
\end{align}
We denote this expanded variance as the marginal variance, $V_{\text{marg}}$.

Since the conditional mean is linear in $\theta_C$ and $\pi_D$ is Normal, the marginal sampling distribution is again Normal:
\begin{equation} \label{eq:app_marginal_distribution}
    \hat{\delta} \mid \delta \sim \mathcal{N}\Big( (1 - W)(\mu_{A,C} - \mu_{D,C}) + \delta, \; V_{\text{marg}} \Big).
\end{equation}

\paragraph{Closed-Form Solutions}
Evaluating the probability that this marginal distribution crosses $k$ yields the average power:
\begin{equation} \label{eq:app_avgpower_closed_form}
    \avgPow(\delta) = 1 - \Phi\left( \frac{q + z_p \sqrt{V_{\delta}} - \big[(1 - W)(\mu_{A,C} - \mu_{D,C}) + \delta\big]}{\sqrt{V_{\text{marg}}}} \right).
\end{equation}
Setting $\delta = 0$ yields the average T1E rate: 
\begin{equation} \label{eq:app_avgt1e_closed_form}
    \avgTIE = 1 - \Phi\left( \frac{q + z_p \sqrt{V_{\delta}} - (1 - W)(\mu_{A,C} - \mu_{D,C})}{\sqrt{V_{\text{marg}}}} \right).
\end{equation}

with $V_{\delta} = \frac{\sigma_{A,C}^2 v_C^2}{\sigma_{A,C}^2 + v_C^2} + v_T^2$.

\section{Mathematical properties of the average metrics}

\paragraph{Property 1: Predictive metrics closer to 50\%}
Evaluating the conditional metric at the center of the design prior ($\theta_C^* = \mu_{D,C}$), the numerators for the $Z$-scores of conditional power and average power are identical. However, as proven in Equation~\eqref{eq:app_marginal_variance}, the marginal variance strictly exceeds the sampling variance ($V_{\text{marg}} > V_{\text{samp}}$). 
Because the denominator of $Z_{\text{marg}}$ is larger, its absolute magnitude must be smaller:
$$
    |Z_{\text{marg}}| < |Z_{\text{cond}}| \quad \text{whenever the numerator is nonzero}.
$$
Since $\Phi(Z)$ is monotonically increasing and centered at $\Phi(0) = 0.5$, dividing by a larger variance shrinks the $Z$-score toward zero, forcing the average metric to be closer to $0.5$ compared to its conditional counterpart. 

\paragraph{Property 2: Exact Control under Alignment}
Assume the design prior exactly matches the analysis prior ($\mu_{D,C} = \mu_{A,C}$ and $\sigma_{D,C}^2 = \sigma_{A,C}^2$). 
The numerator bias term becomes $(1 - W)(\mu_{A,C} - \mu_{A,C}) = 0$. 
In the denominator, substituting $\sigma_{A,C}^2$ into $V_{\text{marg}}$ and expanding weights $W$:
\begin{align*}
    W^2 v_C^2 + (1 - W)^2 \sigma_{A,C}^2 &= \left( \frac{\sigma_{A,C}^2}{\sigma_{A,C}^2 + v_C^2} \right)^2 v_C^2 + \left( \frac{v_C^2}{\sigma_{A,C}^2 + v_C^2} \right)^2 \sigma_{A,C}^2 \\
    &= \frac{\sigma_{A,C}^2 v_C^2 (\sigma_{A,C}^2 + v_C^2)}{(\sigma_{A,C}^2 + v_C^2)^2} = \frac{\sigma_{A,C}^2 v_C^2}{\sigma_{A,C}^2 + v_C^2} = \sigma_{\text{post},C}^2.
\end{align*}
Therefore, $V_{\text{marg}} = \sigma_{\text{post},C}^2 + v_T^2$, which is exactly $V_{\delta}$. Substituting these back into Equation~\eqref{eq:app_avgt1e_closed_form}:
$$
    \avgTIE = 1 - \Phi\left( \frac{q + z_p \sqrt{V_{\delta}}}{\sqrt{V_{\delta}}} \right) = 1 - \Phi\left( \frac{q}{\sqrt{V_{\delta}}} + z_p \right).
$$
If $q = 0$, $\avgTIE = 1 - \Phi(z_p)$. Because $z_p = \Phi^{-1}(p)$ and the nominal significance level is $\alpha = 1 - p$, we achieve exact control: $\avgTIE = \alpha$. Note that this exact control holds at $q = 0$; for a non-zero reference value $q \neq 0$ (e.g., a non-inferiority margin), the result is $\avgTIE = 1 - \Phi(q/\sqrt{V_{\delta}} + z_p) \neq \alpha$.

\paragraph{Property 3: The Triviality of a vague Design Prior}
Assuming a positive weight assigned to the historical prior ($(1 - W) > 0$), we evaluate the limit of the marginal variance as design uncertainty approaches infinity ($\sigma_{D,C}^2 \to \infty$):
$$
    \lim_{\sigma_{D,C}^2 \to \infty} V_{\text{marg}} = \lim_{\sigma_{D,C}^2 \to \infty} \Big[ W^2 v_C^2 + v_T^2 + (1 - W)^2 \sigma_{D,C}^2 \Big] = \infty.
$$
Substituting $V_{\text{marg}} \to \infty$ into Equation~\eqref{eq:app_avgt1e_closed_form}:
$$
    \lim_{\sigma_{D,C}^2 \to \infty} \avgTIE = 1 - \Phi\left( \lim_{\sigma_{D,C}^2 \to \infty} \frac{q + z_p \sqrt{V_{\delta}} - (1 - W)(\mu_{A,C} - \mu_{D,C})}{\sqrt{V_{\text{marg}}}} \right) = 1 - \Phi(0) = 0.5.
$$
Notably, the average power also converges to exactly 0.5.






%%% Local Variables: 
%%% mode: latex
%%% TeX-master: "../MasterThesisSfS"
%%% End: 
